\documentclass[fleqn,twoside,twocolumn,nofootinbib,showkeys]{revtex4}
\usepackage[sec,doi,cyr]{ujp_UTF8}
%\usepackage{graphicx}% Include figure files
%\usepackage{dcolumn}
\usepackage{color}
%\usepackage{color,soul}
%\usepackage{url}
%\usepackage{makecell}
%\draft % marks overfull lines with a black rule on the right

\begin{document}

Шановні Редакторе та Рецензенте!

Щиро дякуємо за те, що ви знайшли час для рецензування нашої статті
``Застосування моделей машинного навчання для ефективної оцінки рухливості носіїв у кремнії''.
Ваші конструктивні пропозиції допомогли  поліпшити якість  роботи.
Ми ретельно розглянули всі коментарі та внесли відповідні зміни до рукопису.
Нижче ми надаємо детальні відповіді на кожен коментар.
Сподіваємося, що оновлений рукопис краще відповідає вашим очікуванням та стандартам для публікації в ``Українському фізичному журналі''.



\vspace{1cm}
\noindent
\textcolor[rgb]{0.00,0.50,1.00}{\textbf{Зауваження~1.}}
\emph{У вступі не розкрито стан методів машинного навчання взагалі, та використання цих методів у дослідженні напівпровідникових властивостей, зокрема. Надалі, у розділі 2.2, надано більше інформації про машинне навчання, проте, про це має бути мова і у вступі.}

\noindent
\textcolor[rgb]{0.51,0.00,0.00}{\textbf{Відповідь:}}
Вступ доповнений інформацією щодо використання методів машинного навчання при розгляді рухливості носіїв заряду в напівпровідниках
(сторінка 3, останній абзац лівої колонки; перший абзац правої).
Питання використання цих методів у дослідженні напівпровідникових властивостей загалом є набагато більш ширшим та добре висвітлено у низці оглядів.
На деякі з них надано посилання в роботі.


\vspace{1cm}
\noindent
\textcolor[rgb]{0.00,0.50,1.00}{\textbf{Зауваження~2.}}
\emph{У вступі достатньо ретельно описано математичне та фізичне підгрунтя проблеми рухливості носіїв струму, проте не завжди надано відповідні посилання.
Ті посилання, що надані, достатньо застарілі, наприклад, з 10 перших посилань тільки 2 посилання 2004 та 2010 років і жодного посилання в діапазоні 10 років.}

\noindent
\textcolor[rgb]{0.51,0.00,0.00}{\textbf{Відповідь:}}
Вступ доповнено посиланнями, зокрема на недавні джерела.
Зазначимо, що більшість робіт, де були запропоновані загальновживані моделі для опису рухливості,
опубліковані досить давно; 
прагнення процитувати саме оригінальні статті і було основною причиною певної застарілості.
Проте ми абсолютно згодні із зауваженням, певна частина посилань оновлена.
З 10 перших посилань 6 відповідають діапазону останніх 10 років.


\vspace{1cm}
\noindent
\textcolor[rgb]{0.00,0.50,1.00}{\textbf{Зауваження~3.}}
\emph{Як на мене, серед посилань, наведених в цьому рукописі, тільки 9 не старші за 10 років (9 з 23, що складає приблизно 40 \%), а мають бути не менше 70.}

\noindent
\textcolor[rgb]{0.51,0.00,0.00}{\textbf{Відповідь:}}
Після модифікації роботи відповідно до перших двох зауважень, посилання не старші за 10 років складають біля 72 \%.


\vspace{1cm}
З повагою, 

Олег Оліх

\end{document}

